\section{การติดตั้งเฟิร์มแวร์และการตั้งค่าซอฟต์แวร์}
\label{sec:firmware}

เฟิร์มแวร์ของเครื่องสเปกโทรโฟโตมิเตอร์พัฒนาขึ้นบนระบบนิเวศ Arduino โดยปรับแต่งให้ทำงานบนสถาปัตยกรรม ARM Cortex-M4F ของไมโครคอนโทรลเลอร์ ATSAMD51 บนบอร์ด Seeed Studio Wio Terminal พร้อมระบบบริหารจัดการหน่วยความจำและจอแสดงผลที่มีประสิทธิภาพสูง

\subsection{ไลบรารีที่จำเป็นสำหรับระบบ}
การพัฒนาและคอมไพล์โปรแกรมต้องการไลบรารีมาตรฐานดังต่อไปนี้

\begin{enumerate}
    \item \textbf{Seeed\_Arduino\_LCD} ไลบรารีควบคุมจอแสดงผล TFT LCD ความละเอียด 320x240 พิกเซลผ่านบัส SPI3 เฉพาะตัวของ Wio Terminal
    \item \textbf{Adafruit\_TCS34725} ไลบรารีสื่อสารกับเซนเซอร์วัดค่าสีแบบ I2C ควบคุมเวลาการประมวลผล (Integration Time) และเกนขยายสัญญาณ
    \item \textbf{Adafruit\_NeoPixel} ไลบรารีขับสัญญาณดิจิทัลควบคุมโมดูลหลอดไฟ RGB WS2812B
    \item \textbf{Seeed\_Arduino\_FS และ Seeed\_SD} ไลบรารีจัดการระบบไฟล์และบันทึกข้อมูลลงในการ์ดหน่วยความจำ microSD Card
\end{enumerate}

\subsection{ข้อควรระวังสำคัญด้านไลบรารีจอแสดงผล}
\begin{academicbox}{การป้องกันปัญหาจอขาวและการตั้งค่าไลบรารี}
    ในการคอมไพล์โปรแกรมสำหรับบอร์ด Wio Terminal จะต้องใช้ไลบรารี \textbf{Seeed\_Arduino\_LCD} ที่ติดตั้งมาพร้อมกับ Core Seeeduino SAMD เท่านั้น โดยห้ามใช้ไลบรารี TFT\_eSPI ทั่วไปของ Bodmer เนื่องจากมีการกำหนดขาสัญญาณ SPI ผิดพลาดและส่งผลให้จอแสดงผลเป็นสีขาวสว่างค้าง
    
    นอกจากนี้ เพื่อให้ระบบสามารถแสดงผลฟอนต์ภาษาไทย THSarabunPSK ได้อย่างคมชัด จะต้องเปิดการทำงานของคุณสมบัติ Smooth Font ในไฟล์ \texttt{User\_Setup.h} ของ Seeed\_Arduino\_LCD โดยเปิดใช้งานบรรทัด:
    \begin{center}
        \texttt{\#define SMOOTH\_FONT}
    \end{center}
\end{academicbox}

\subsection{คำสั่งคอมไพล์และอัปโหลดโปรแกรม}
ผู้ใช้งานสามารถคอมไพล์และอัปโหลดโปรแกรมผ่านโปรแกรม Arduino CLI บนระบบปฏิบัติการ macOS หรือ Linux โดยใช้ชุดคำสั่งดังต่อไปนี้

\begin{center}
\begin{minipage}{0.92\textwidth}
\begin{verbatim}
# Compile sketch for Wio Terminal
arduino-cli compile --fqbn Seeeduino:samd:seeed_wio_terminal firmware

# Upload sketch via USB port
arduino-cli upload -p /dev/cu.usbmodem2101 \
    --fqbn Seeeduino:samd:seeed_wio_terminal firmware
\end{verbatim}
\end{minipage}
\end{center}

กรณีที่เกิดปัญหาพอร์ต USB หลุดการเชื่อมต่อ สามารถแก้ไขได้โดยการเลื่อนสวิตช์เปิดปิดด้านข้างตัวเครื่องลง 2 ครั้งอย่างรวดเร็ว (Double-click Touch Reset) เพื่อเข้าสู่โหมด Bootloader จากนั้นคัดลอกไฟล์ไบนารี \texttt{spectrometer\_npk2026.uf2} ไปวางในไดรฟ์ USB ที่ปรากฏขึ้นได้อย่างสะดวก
